%% ╔══════════════════════════════════════════════════════════════════╗
%% ║         CORIOTLAB — PLANTILLA PRESENTACIÓN ACADÉMICA            ║
%% ║         LaTeX Beamer · diseño basado en Brand Book CORIOTLAB    ║
%% ╚══════════════════════════════════════════════════════════════════╝
%%
%% ═══════════════════════════════════════════════════════════════════
%% FICHA DE USO
%% ═══════════════════════════════════════════════════════════════════
%% PARA QUÉ: Presentaciones académicas y profesionales de proyectos,
%%           resultados de investigación o ponencias de CORIOTLAB.
%% PARA QUIÉN: Presentaciones internas (reuniones, avances) y externas
%%             (conferencias, clientes, evaluaciones).
%% ESTRUCTURA: 8 diapositivas tipo — seleccionar y adaptar las necesarias.
%%   1. Portada          5. Dos columnas
%%   2. Agenda           6. Código
%%   3. Separador        7. Resultados
%%   4. Contenido        8. Cierre
%% COMPILAR: xelatex -interaction=nonstopmode plantilla_presentacion.tex
%%           (doble pasada — ejecutar el comando dos veces)
%% OUTPUT:   plantilla_presentacion.pdf (en la misma carpeta)
%% NOTA: NO usa membrete PNG — diseño propio con TikZ basado en Brand Book.
%% ═══════════════════════════════════════════════════════════════════

%% ═══════════════════════════════════════════════════════════════════
%% ZONA DE CONFIGURACIÓN — editar aquí antes de compilar
%% ═══════════════════════════════════════════════════════════════════

%% TÍTULO DE LA PRESENTACIÓN: Nombre del proyecto o tema a exponer.
%%   Aparece en portada, cierre y en el pie de cada diapositiva.
%%   Máximo 60 caracteres para que no se trunque en el footline.
%%   Ejemplo: "Sistema de Monitoreo IoT Distribuido"
\newcommand{\PresTitle}{[Título de la presentación]}

%% SUBTÍTULO: Complemento técnico o contexto del título.
%%   Aparece debajo del título en la portada.
%%   Ejemplo: "Red de sensores ESP32 con MQTT, InfluxDB y Grafana"
\newcommand{\PresSubtitle}{[Subtítulo: contexto o enfoque técnico]}

%% AUTOR: Nombre completo del presentador.
\newcommand{\PresAutor}{[Nombre completo del presentador]}

%% CARGO: Rol del presentador en CORIOTLAB.
%%   Ejemplo: "Investigador --- CORIOTLAB"
\newcommand{\PresCargo}{[Cargo --- CORIOTLAB]}

%% PROYECTO: Nombre del proyecto al que pertenece la presentación.
%%   Ejemplo: "Proyecto de Monitoreo Ambiental"
\newcommand{\PresProyecto}{[Nombre del proyecto]}

%% FECHA: Fecha de la presentación.
%%   Formato: DD de mes de AAAA — Ejemplo: 30 de junio de 2025
\newcommand{\PresFecha}{[DD de mes de AAAA]}

%% INSTITUCIÓN: Entidad que respalda la presentación.
%%   Ejemplo: "CORIOTLAB --- ITM"
\newcommand{\PresInstitucion}{CORIOTLAB --- ITM}

%% ═══════════════════════════════════════════════════════════════════

\documentclass[aspectratio=169, 12pt]{beamer}

%% --- Lenguaje
\usepackage[spanish, es-nodecimaldot, es-noshorthands]{babel}

%% --- Fuentes
\usepackage{fontspec}
\setmainfont{Inter}
\newfontfamily\FuenteTitulo{MuseoModerno}
\newfontfamily\FuenteCodigo{Space Mono}

%% --- Color
\usepackage{xcolor}
\definecolor{AzulITM}     {HTML}{102D69}
\definecolor{Magenta}     {HTML}{C14894}
\definecolor{AzulDigital} {HTML}{56ACDE}
\definecolor{GrisPizarra} {HTML}{2F2F2F}
\definecolor{GrisClaro}   {HTML}{F2F2F2}
\definecolor{GrisMedio}   {HTML}{AAAAAA}
\definecolor{GrisLinea}   {HTML}{DDDDDD}
\definecolor{Blanco}      {HTML}{FFFFFF}

%% --- Gráficos y TikZ
\usepackage{tikz}
\usepackage{graphicx}

%% --- Cajas
\usepackage[most]{tcolorbox}

%% --- Código (frames con lstlisting DEBEN llevar [fragile])
\usepackage{listings}

%% --- Tablas
\usepackage{array}
\usepackage{colortbl}

%% ─────────────────────────────────────────────────────────────────
%% TEMA BEAMER
%% ─────────────────────────────────────────────────────────────────
\usetheme{default}
\useinnertheme{default}
\useoutertheme{default}

\setbeamercolor{normal text}     {fg=GrisPizarra, bg=white}
\setbeamercolor{structure}       {fg=AzulITM}
\setbeamercolor{palette primary} {bg=AzulITM, fg=white}
\setbeamercolor{frametitle}      {bg=AzulITM, fg=white}
\setbeamercolor{title}           {fg=white}
\setbeamercolor{subtitle}        {fg=AzulDigital}
\setbeamercolor{block title}     {bg=AzulITM, fg=white}
\setbeamercolor{block body}      {bg=GrisClaro, fg=GrisPizarra}
\setbeamercolor{alerted text}    {fg=Magenta}
\setbeamercolor{item}            {fg=AzulITM}
\setbeamercolor{footline}        {fg=GrisMedio, bg=white}
\setbeamercolor{enumerate item}  {fg=AzulITM}

\setbeamerfont{title}      {family=\FuenteTitulo, size=\huge,      series=\bfseries}
\setbeamerfont{subtitle}   {family=\FuenteTitulo, size=\large,     series=\mdseries}
\setbeamerfont{frametitle} {family=\FuenteTitulo, size=\large,     series=\bfseries}
\setbeamerfont{block title}{family=\FuenteTitulo, size=\normalsize, series=\bfseries}
\setbeamerfont{footline}   {size=\scriptsize}

%% Bullets Beamer nativos con > (NO usar enumitem en Beamer)
\setbeamertemplate{itemize item}   {\color{AzulITM}\textbf{>}}
\setbeamertemplate{itemize subitem}{\color{AzulDigital}\small\textbf{>}}
\setbeamertemplate{enumerate item} {\color{AzulITM}\textbf{\insertenumlabel.}}
\setbeamertemplate{navigation symbols}{}

\setbeamertemplate{frametitle}{%
  \vspace{0.3em}\insertframetitle\vspace{0.2em}}

%% ─────────────────────────────────────────────────────────────────
%% FOOTLINE
%% ─────────────────────────────────────────────────────────────────
\setbeamertemplate{footline}{%
  \leavevmode%
  \begin{beamercolorbox}[wd=\paperwidth, ht=0.6pt, dp=0pt]{structure}%
  \end{beamercolorbox}%
  \begin{beamercolorbox}[wd=\paperwidth, ht=2.0em, dp=0.8em,
      leftskip=1em, rightskip=1em]{footline}%
    {\FuenteTitulo\tiny\color{AzulITM}\textbf{< CORIOTLAB}}%
    \hfill%
    {\tiny\color{GrisMedio}\PresTitle}%
    \hfill%
    {\tiny\color{GrisMedio}\insertframenumber\,/\,\inserttotalframenumber}%
  \end{beamercolorbox}%
}
\setbeamertemplate{headline}{}

%% ─────────────────────────────────────────────────────────────────
%% CÓDIGO lstlisting
%% ─────────────────────────────────────────────────────────────────
\lstdefinestyle{coriotcode}{
  basicstyle=\FuenteCodigo\scriptsize,
  backgroundcolor=\color{GrisClaro},
  commentstyle=\color{GrisMedio}\itshape,
  keywordstyle=\color{AzulITM}\bfseries,
  stringstyle=\color{Magenta},
  numberstyle=\tiny\color{GrisMedio},
  numbers=left, numbersep=8pt,
  breaklines=true, frame=none,
  tabsize=2, showstringspaces=false
}
\lstset{style=coriotcode}

%% ─────────────────────────────────────────────────────────────────
%% CAJAS tcolorbox
%% ─────────────────────────────────────────────────────────────────
\tcbset{
  cajaNota/.style={
    colback=AzulDigital!10, colframe=AzulDigital,
    fonttitle=\FuenteTitulo\small\bfseries, coltitle=white,
    boxrule=0pt, leftrule=4pt, title=Nota,
    left=6pt, top=4pt, bottom=4pt
  },
  cajaCode/.style={
    colback=GrisClaro, colframe=AzulDigital,
    boxrule=0pt, leftrule=4pt,
    left=6pt, top=4pt, bottom=4pt
  },
  cajaIndicador/.style={
    colback=GrisClaro, colframe=GrisLinea,
    boxrule=0.5pt, halign=center,
    left=4pt, right=4pt, top=6pt, bottom=6pt
  }
}

%% ─────────────────────────────────────────────────────────────────
%% MACRO: LOGO CORIOTLAB en TikZ
%% ─────────────────────────────────────────────────────────────────
\newcommand{\LogoCoriotTikZ}[2]{%
  \begin{tikzpicture}[baseline=-0.3ex]
    \node[anchor=base east,  text=#1, font=\FuenteTitulo\Large\bfseries]
      at (0,0)     {<};
    \node[anchor=base west,  text=#2, font=\FuenteTitulo\Large\bfseries]
      at (0.05cm,0){CORIOTLAB};
    \fill[AzulDigital] (-0.85cm, 0.72cm) rectangle (-0.65cm, 0.52cm);
    \fill[Magenta]     (-0.65cm, 0.48cm) rectangle (-0.48cm, 0.32cm);
    \fill[AzulITM]     (-0.85cm, 0.28cm) rectangle (-0.72cm, 0.16cm);
  \end{tikzpicture}%
}

%% =================================================================
\begin{document}

%% -----------------------------------------------------------------
%% DIAPOSITIVA 1: PORTADA [plain — sin footline]
%% Esta diapositiva se genera automáticamente con los datos de la
%% Zona de Configuración. No modificar la estructura.
%% -----------------------------------------------------------------
\begin{frame}[plain]
  \begin{tikzpicture}[remember picture, overlay]
    \fill[AzulITM] (current page.south west) rectangle (current page.north east);
    %% Píxeles decorativos esquina superior derecha
    \fill[AzulDigital]
      ([xshift=-2.4cm, yshift=-1.0cm] current page.north east)
      rectangle ([xshift=-1.9cm, yshift=-1.5cm] current page.north east);
    \fill[Magenta]
      ([xshift=-1.8cm, yshift=-1.5cm] current page.north east)
      rectangle ([xshift=-1.4cm, yshift=-1.9cm] current page.north east);
    \fill[AzulDigital!60]
      ([xshift=-2.4cm, yshift=-2.0cm] current page.north east)
      rectangle ([xshift=-2.0cm, yshift=-2.4cm] current page.north east);
    \fill[Magenta!70]
      ([xshift=-1.3cm, yshift=-1.0cm] current page.north east)
      rectangle ([xshift=-1.0cm, yshift=-1.3cm] current page.north east);
    %% Franja inferior AzulDigital
    \fill[AzulDigital]
      (current page.south west) rectangle ([yshift=0.5cm] current page.south east);
  \end{tikzpicture}

  \vspace{1.2cm}
  \begin{center}
    \LogoCoriotTikZ{AzulDigital}{white}\\[1.2em]
    {\FuenteTitulo\huge\color{white}\textbf{\PresTitle}}\\[0.5em]
    {\color{AzulDigital}\rule{0.6\textwidth}{1.5pt}}\\[0.5em]
    {\FuenteTitulo\large\color{AzulDigital}\PresSubtitle}\\[1.4em]
    {\small\color{white}\PresAutor\ \ \textbullet\ \ \PresCargo}\\[0.2em]
    {\small\color{white}\PresInstitucion\ \ \textbullet\ \ \PresFecha}
  \end{center}
\end{frame}

%% -----------------------------------------------------------------
%% DIAPOSITIVA 2: AGENDA
%% Liste los temas o secciones principales de la presentación.
%% Entre 3 y 6 ítems recomendados para una buena legibilidad.
%% Los números se asignan automáticamente por el template enumerate.
%% -----------------------------------------------------------------
\begin{frame}{Agenda}
  \begin{tikzpicture}[overlay, remember picture]
    %% Barra lateral AzulITM (acento visual izquierdo)
    \fill[AzulITM]
      ([xshift=0.6cm,  yshift=-0.85cm] current page.north west) rectangle
      ([xshift=0.86cm, yshift=1.1cm]   current page.south west);
  \end{tikzpicture}

  \hspace{0.7cm}
  \begin{minipage}{0.88\textwidth}
    %% INSTRUCCIONES: Reemplazar cada ítem con los temas de su presentación.
    %% NO agregar [label=...] — usar el template nativo de Beamer.
    \begin{enumerate}
      \item {[}Primer tema o sección de la presentación{]}
      \item {[}Segundo tema o sección{]}
      \item {[}Tercer tema o sección{]}
      \item {[}Cuarto tema o sección{]}
      \item {[}Quinto tema: conclusiones y trabajo futuro{]}
    \end{enumerate}
  \end{minipage}
\end{frame}

%% -----------------------------------------------------------------
%% DIAPOSITIVA 3: SEPARADOR DE SECCIÓN [plain — sin footline]
%% Usar al inicio de cada sección temática principal.
%% Duplicar este frame para cada sección nueva de la presentación.
%% -----------------------------------------------------------------
\begin{frame}[plain]
  \begin{tikzpicture}[remember picture, overlay]
    \fill[AzulITM] (current page.south west) rectangle (current page.north east);
    %% Símbolo "<" decorativo al 15% opacidad — no modificar
    \node[opacity=0.15, anchor=center] at (current page.center)
      {\scalebox{16}{\color{white}\FuenteTitulo<}};
    %% NÚMERO DE SECCIÓN: cambiar "01" por el número correspondiente
    \node[anchor=south west] at ([xshift=2.2cm, yshift=1.2cm] current page.south west)
      {\FuenteTitulo\Huge\color{AzulDigital}\textbf{01}};
    %% TÍTULO DE SECCIÓN: reemplazar con el nombre de la sección
    \node[anchor=south west] at ([xshift=2.2cm, yshift=3.2cm] current page.south west)
      {\FuenteTitulo\Large\color{white}\textbf{[Título de la Sección]}};
  \end{tikzpicture}
\end{frame}

%% -----------------------------------------------------------------
%% DIAPOSITIVA 4: CONTENIDO ESTÁNDAR
%% Diapositiva de texto con viñetas. Estructura más común.
%% Incluye barra lateral AzulITM y caja de nota opcional.
%% -----------------------------------------------------------------
\begin{frame}{[Título de la diapositiva]}
  \begin{tikzpicture}[overlay, remember picture]
    \fill[AzulITM]
      ([xshift=0.6cm,  yshift=-0.85cm] current page.north west) rectangle
      ([xshift=0.85cm, yshift=1.1cm]   current page.south west);
  \end{tikzpicture}

  \hspace{0.55cm}
  \begin{minipage}{0.88\textwidth}
    %% INSTRUCCIONES: Reemplazar el texto de introducción y los ítems.
    %% NO agregar [label=...] — usar el template nativo de Beamer.
    {[}Frase introductoria opcional que contextualice los puntos siguientes.{]}
    \vspace{0.4em}
    \begin{itemize}
      \item \textbf{[Punto clave 1:]} {[}desarrollo del punto en 1--2 líneas.{]}
      \item \textbf{[Punto clave 2:]} {[}desarrollo del punto en 1--2 líneas.{]}
      \item \textbf{[Punto clave 3:]} {[}desarrollo del punto en 1--2 líneas.{]}
    \end{itemize}

    %% NOTA OPCIONAL: eliminar el bloque cajaNota si no es necesario.
    \vspace{0.5em}
    \begin{tcolorbox}[cajaNota, top=3pt, bottom=3pt]
      {[}Nota destacada: información clave que el audiencia debe retener.{]}
    \end{tcolorbox}
  \end{minipage}
\end{frame}

%% -----------------------------------------------------------------
%% DIAPOSITIVA 5: DOS COLUMNAS
%% Para comparativas, listas duales o texto + imagen/diagrama.
%% Ajuste el ancho de columnas (0.52 / 0.44) según contenido.
%% -----------------------------------------------------------------
\begin{frame}{[Título de la diapositiva de dos columnas]}
  \begin{columns}[T]
    \begin{column}{0.52\textwidth}
      %% COLUMNA IZQUIERDA: texto, listas o comparativas
      \textbf{\color{AzulITM}{[}Encabezado columna izquierda{]}:}
      \begin{itemize}
        \item {[}Ítem 1 de la lista izquierda.{]}
        \item {[}Ítem 2 de la lista izquierda.{]}
        \item {[}Ítem 3 de la lista izquierda.{]}
      \end{itemize}
      \vspace{0.4em}
      \textbf{\color{AzulITM}{[}Segundo encabezado (opcional){]}:}
      \begin{itemize}
        \item {[}Ítem adicional 1.{]}
        \item {[}Ítem adicional 2.{]}
      \end{itemize}
    \end{column}
    \begin{column}{0.44\textwidth}
      %% COLUMNA DERECHA: imagen, diagrama o texto complementario.
      %% Para insertar imagen: reemplazar el fcolorbox por:
      %%   \includegraphics[width=\textwidth]{nombre-imagen.png}
      \begin{center}
        \fcolorbox{GrisLinea}{GrisClaro}{%
          \begin{minipage}[c][5.2cm][c]{0.9\textwidth}
            \centering
            {\color{GrisMedio}\small {[}Diagrama / imagen / figura{]}\\[0.4em]
            {[}descripción o nombre del archivo{]}\\[0.4em]
            \FuenteCodigo\tiny nombre-archivo.png}
          \end{minipage}%
        }
      \end{center}
    \end{column}
  \end{columns}
\end{frame}

%% -----------------------------------------------------------------
%% DIAPOSITIVA 6: CÓDIGO [fragile — OBLIGATORIO con lstlisting]
%% Para mostrar fragmentos de código fuente.
%% Lenguajes disponibles: Python, C++, bash, C, Java, etc.
%% -----------------------------------------------------------------
\begin{frame}[fragile]{[Título: descripción del código mostrado]}
  \begin{tcolorbox}[cajaCode, top=2pt, bottom=2pt]
%% INSTRUCCIONES: Reemplazar el código de ejemplo con el código real.
%% Cambiar "language=Python" por el lenguaje correspondiente.
%% Eliminar la caption si no es necesaria.
\begin{lstlisting}[language=Python]
# [Reemplazar con código real del proyecto]
def nombre_funcion(parametro):
    """[Descripción breve de la función.]"""
    resultado = [logica-aqui]
    return resultado

# [Llamada de ejemplo]
salida = nombre_funcion([argumento])
print(salida)
\end{lstlisting}
  \end{tcolorbox}
  {\small\color{GrisMedio}\itshape
    {[}Descripción breve del código: qué hace, por qué es relevante,
    qué resultado produce.{]}}
\end{frame}

%% -----------------------------------------------------------------
%% DIAPOSITIVA 7: RESULTADOS
%% Para presentar métricas e indicadores de desempeño.
%% Ajuste los valores en cajaIndicador y la tabla según sus datos.
%% -----------------------------------------------------------------
\begin{frame}{Resultados de Desempeño}
  %% INDICADORES: 3 métricas clave en cajas destacadas.
  %% Cambiar el valor, la unidad y la etiqueta de cada indicador.
  \begin{columns}[T]
    \begin{column}{0.32\textwidth}
      \begin{tcolorbox}[cajaIndicador]
        {\FuenteTitulo\LARGE\color{AzulITM}\textbf{[Valor 1]}}\\[4pt]
        {\small\color{GrisMedio}{[}Indicador 1{]}}
      \end{tcolorbox}
    \end{column}
    \begin{column}{0.32\textwidth}
      \begin{tcolorbox}[cajaIndicador]
        {\FuenteTitulo\LARGE\color{AzulITM}\textbf{[Valor 2]}}\\[4pt]
        {\small\color{GrisMedio}{[}Indicador 2{]}}
      \end{tcolorbox}
    \end{column}
    \begin{column}{0.32\textwidth}
      \begin{tcolorbox}[cajaIndicador]
        {\FuenteTitulo\LARGE\color{AzulITM}\textbf{[Valor 3]}}\\[4pt]
        {\small\color{GrisMedio}{[}Indicador 3{]}}
      \end{tcolorbox}
    \end{column}
  \end{columns}

  %% TABLA DE RESULTADOS: detallar indicadores vs. metas.
  %% Agregar o eliminar filas según la cantidad de indicadores.
  \vspace{0.6em}
  {\small
  \rowcolors{2}{GrisClaro}{white}
  \begin{tabular}{p{3.8cm}
      >{\centering\arraybackslash}p{1.7cm}
      >{\centering\arraybackslash}p{1.7cm}
      >{\centering\arraybackslash}p{1.7cm}
      p{3.2cm}}
    \rowcolor{AzulITM}
    \color{Blanco}\FuenteTitulo\scriptsize\bfseries Indicador &
    \color{Blanco}\FuenteTitulo\scriptsize\bfseries Meta &
    \color{Blanco}\FuenteTitulo\scriptsize\bfseries Logrado &
    \color{Blanco}\FuenteTitulo\scriptsize\bfseries \% Cumpl. &
    \color{Blanco}\FuenteTitulo\scriptsize\bfseries Observación \\
    {[}Indicador 1{]} & {[}Meta{]} & {[}Logrado{]} & {[}XX\%{]} & {[}Nota breve.{]} \\
    {[}Indicador 2{]} & {[}Meta{]} & {[}Logrado{]} & {[}XX\%{]} & {[}Nota breve.{]} \\
    {[}Indicador 3{]} & {[}Meta{]} & {[}Logrado{]} & {[}XX\%{]} & {[}Nota breve.{]} \\
  \end{tabular}
  }
\end{frame}

%% -----------------------------------------------------------------
%% DIAPOSITIVA 8: CIERRE [plain — sin footline]
%% Se genera automáticamente con datos de la Zona de Configuración.
%% No modificar la estructura. El lector recordará: logo + gracias +
%% slogan + contacto.
%% -----------------------------------------------------------------
\begin{frame}[plain]
  \begin{tikzpicture}[remember picture, overlay]
    \fill[AzulITM] (current page.south west) rectangle (current page.north east);
    \fill[AzulDigital]
      (current page.south west) rectangle ([yshift=0.5cm] current page.south east);
    \fill[AzulDigital]
      ([xshift=-2.4cm, yshift=-1.0cm] current page.north east)
      rectangle ([xshift=-1.9cm, yshift=-1.5cm] current page.north east);
    \fill[Magenta]
      ([xshift=-1.8cm, yshift=-1.5cm] current page.north east)
      rectangle ([xshift=-1.4cm, yshift=-1.9cm] current page.north east);
  \end{tikzpicture}

  \begin{center}
    \vspace{1cm}
    \LogoCoriotTikZ{AzulDigital}{white}\\[1.2em]
    {\FuenteTitulo\Huge\color{white}\textbf{¡Gracias!}}\\[0.4em]
    {\color{AzulDigital}\rule{0.4\textwidth}{1.5pt}}\\[0.5em]
    {\FuenteCodigo\normalsize\color{AzulDigital}%
      CONTROL\ \ \textbullet\ \ ROB\'OTICA\ \ \textbullet\ \ IOT}\\[1.2em]
    {\small\color{white}www.coriotlab.co\ \ |\ \ info@coriotlab.co}
  \end{center}
\end{frame}

\end{document}
